\documentclass[11pt]{article}
\usepackage[T1]{fontenc}
\usepackage[ngerman]{babel}
\usepackage{amsmath,amssymb,amsthm}
\usepackage{geometry}
\geometry{margin=2.5cm}

\title{Quaternionische Formulierung der Rotationsdynamik \\ 
	und eine EABC-Interpretation (Bamberg-Modell)}
\author{}
\date{}

\begin{document}
	
	\maketitle
	
	\begin{abstract}
		Wir formulieren die klassische Rotationsdynamik vollständig in der Sprache der Quaternionen und zeigen, dass sich daraus eine natürliche vierkanalige Struktur ergibt. Diese Struktur erlaubt eine Interpretation im Rahmen eines EABC-Schemas, in dem der skalare Anteil als Energie- bzw. Normkanal fungiert, während die drei imaginären Komponenten gerichtete Dynamik repräsentieren. Die Arbeit liefert eine kompakte algebraische Darstellung der Rotationsmechanik und schlägt eine Brücke zu arithmetischen Modellen auf Basis von Vier-Quadrate-Zerlegungen.
	\end{abstract}
	
	\section{Quaternionen und Rotationen}
	
	Ein Quaternion ist definiert als
	\begin{equation}
		q = q_0 + q_1 i + q_2 j + q_3 k,
	\end{equation}
	mit Norm
	\begin{equation}
		\|q\|^2 = q_0^2 + q_1^2 + q_2^2 + q_3^2.
	\end{equation}
	
	Für Rotationen verwendet man Einheitsquaternionen:
	\begin{equation}
		\|q\| = 1.
	\end{equation}
	
	Eine Rotation um die Achse $\mathbf n = (n_x,n_y,n_z)$ mit Winkel $\theta$ wird beschrieben durch
	\begin{equation}
		q(\theta,\mathbf n)
		=
		\cos\frac{\theta}{2}
		+
		(n_x i+n_y j+n_z k)\sin\frac{\theta}{2}.
	\end{equation}
	
	\section{Kinematik in Quaternionenform}
	
	Die klassische Winkelkinematik lautet:
	\begin{equation}
		\theta(t) = \theta_0 + \omega_0 t + \frac12 \alpha t^2.
	\end{equation}
	
	Quaternionisch ergibt sich:
	\begin{equation}
		q(t)
		=
		\cos\frac{\theta(t)}{2}
		+
		\mathbf n \sin\frac{\theta(t)}{2}.
	\end{equation}
	
	Die Winkelgeschwindigkeit wird als rein imaginäres Quaternion geschrieben:
	\begin{equation}
		\Omega = \omega_x i + \omega_y j + \omega_z k.
	\end{equation}
	
	Die Bewegungsgleichung lautet:
	\begin{equation}
		\dot q = \frac12 q \Omega.
	\end{equation}
	
	\section{Winkelbeschleunigung und Dynamik}
	
	Die Winkelbeschleunigung ist
	\begin{equation}
		A_\alpha = \alpha_x i + \alpha_y j + \alpha_z k.
	\end{equation}
	
	Damit gilt:
	\begin{equation}
		\Omega(t) = \Omega_0 + A_\alpha t.
	\end{equation}
	
	\section{Quaternionenprodukt und Vektoranalysis}
	
	Für zwei reine Vektorquaternionen
	\begin{equation}
		a = a_x i + a_y j + a_z k,\quad
		b = b_x i + b_y j + b_z k
	\end{equation}
	gilt:
	\begin{equation}
		ab
		=
		-\mathbf a \cdot \mathbf b
		+
		\mathbf a \times \mathbf b.
	\end{equation}
	
	\section{Drehmoment und Drehimpuls}
	
	Das Drehmoment:
	\begin{equation}
		T = \tau_x i + \tau_y j + \tau_z k,
	\end{equation}
	ist gegeben durch den Vektorteil:
	\begin{equation}
		\boldsymbol\tau = \operatorname{Vec}(rF).
	\end{equation}
	
	Der Drehimpuls:
	\begin{equation}
		\Lambda = L_x i + L_y j + L_z k,
	\end{equation}
	erfüllt im isotropen Fall:
	\begin{equation}
		\Lambda = I \Omega.
	\end{equation}
	
	Im allgemeinen Fall:
	\begin{equation}
		\Lambda = \mathcal I(\Omega),
	\end{equation}
	wobei $\mathcal I$ ein linearer Operator ist.
	
	\section{Energie und Leistung}
	
	Die Rotationsenergie lautet:
	\begin{equation}
		K = \frac12 \|\Omega\|^2 I,
	\end{equation}
	bzw. allgemein:
	\begin{equation}
		K = \frac12 \langle \Omega, \mathcal I(\Omega) \rangle.
	\end{equation}
	
	Die Leistung ergibt sich aus:
	\begin{equation}
		P = \boldsymbol\tau \cdot \boldsymbol\omega.
	\end{equation}
	
	Quaternionisch:
	\begin{equation}
		P = -\operatorname{Scal}(T\Omega).
	\end{equation}
	
	\section{EABC-Interpretation}
	
	Wir identifizieren:
	\begin{equation}
		E \leftrightarrow 1,\quad
		A \leftrightarrow i,\quad
		B \leftrightarrow j,\quad
		C \leftrightarrow k.
	\end{equation}
	
	Ein Zustand ist dann:
	\begin{equation}
		Q = E + Ai + Bj + Ck.
	\end{equation}
	
	Die Norm:
	\begin{equation}
		N(Q) = E^2 + A^2 + B^2 + C^2
	\end{equation}
	ist eine Erhaltungsgröße.
	
	Für reine Rotationen:
	\begin{equation}
		Q =
		\cos\frac{\theta}{2}
		+
		(Ai + Bj + Ck)\sin\frac{\theta}{2},
	\end{equation}
	mit
	\begin{equation}
		A^2 + B^2 + C^2 = 1.
	\end{equation}
	
	\section{Dynamisches Modell}
	
	Ein zeitabhängiger Zustand:
	\begin{equation}
		Q(t) = E(t) + A(t)i + B(t)j + C(t)k
	\end{equation}
	erfüllt:
	\begin{equation}
		\dot Q = \frac12 Q \Omega.
	\end{equation}
	
	Mit:
	\begin{align}
		\Omega &= \omega_A i + \omega_B j + \omega_C k,\\
		\Lambda &= \mathcal I(\Omega),\\
		T &= \dot\Lambda,\\
		K &= \frac12 \langle \Omega,\Lambda\rangle,\\
		P &= \langle T,\Omega\rangle.
	\end{align}
	
	\section{Arithmetische Interpretation}
	
	Sei
	\begin{equation}
		Q_N = E_N + A_N i + B_N j + C_N k
	\end{equation}
	mit
	\begin{equation}
		N = E_N^2 + A_N^2 + B_N^2 + C_N^2.
	\end{equation}
	
	Eine Transformation:
	\begin{equation}
		Q_N \mapsto q Q_N q^{-1}
	\end{equation}
	erhält die Norm und beschreibt eine Rotation im EABC-Raum.
	
	\section{Beispiel: Symmetrische Rotation}
	
	Für die Achse
	\begin{equation}
		\mathbf n = \frac{1}{\sqrt3}(1,1,1)
	\end{equation}
	und $\theta = \frac{2\pi}{3}$ ergibt sich:
	\begin{equation}
		q = \frac{1+i+j+k}{2}.
	\end{equation}
	
	Dies ist ein Hurwitz-Quaternion mit tetraedrischer Symmetrie.
	
	\section{Interpretation}
	
	Die Struktur lässt sich wie folgt lesen:
	\begin{align}
		\text{E} &\rightarrow \text{Skalar, Energie, Norm},\\
		\text{A,B,C} &\rightarrow \text{gerichtete Dynamikachsen}.
	\end{align}
	
	Damit ergibt sich eine klare Trennung:
	\begin{itemize}
		\item Energie als skalare Größe (E-Kanal)
		\item Dynamik als gerichtete Bewegung (A,B,C-Kanäle)
		\item Quaternionenmultiplikation als nichtkommutative Kopplung
	\end{itemize}
	
	\section{Schlussbemerkung}
	
	Die quaternionische Darstellung der Rotationsdynamik liefert nicht nur eine elegante mathematische Formulierung, sondern auch eine natürliche vierdimensionale Struktur. Diese erlaubt eine Interpretation im Sinne eines EABC-Modells, in dem Norm, Energie und gerichtete Dynamik in einer einheitlichen algebraischen Sprache zusammengeführt werden.
	
\end{document}